% =====================================================================
%  CryptoSense AI - Final Year Project Documentation
%  COMSATS University Islamabad, Attock Campus - BS Computer Science
%  Compile with pdfLaTeX (Overleaf: Menu -> Compiler -> pdfLaTeX), run twice
%  so the table of contents and lists are filled in.
% =====================================================================
\documentclass[12pt,a4paper,oneside]{report}

% ---------------------------------------------------------------- packages
\usepackage[a4paper,top=1in,bottom=1in,left=1.5in,right=1in]{geometry}
\usepackage[T1]{fontenc}
\usepackage{newtxtext,newtxmath}          % Times New Roman style text and maths
\usepackage{setspace}
\usepackage{graphicx}
\usepackage{float}
\usepackage{array,tabularx,longtable,booktabs,multirow}
\usepackage[table]{xcolor}
\usepackage{enumitem}
\usepackage{caption}
\usepackage{titlesec}
\usepackage{tocloft}
\usepackage{amsmath}
\usepackage{url}
\usepackage[hidelinks]{hyperref}

\onehalfspacing
\graphicspath{{figures/}}
\setlength{\parindent}{0pt}
\setlength{\parskip}{6pt}

% ------------------------------------------------- project details (EDIT)
\newcommand{\projecttitle}{CryptoSense AI: An Explainable AI-Based\\ Cryptocurrency Market Risk Intelligence Platform}
\newcommand{\projecttitleinline}{CryptoSense AI: An Explainable AI-Based Cryptocurrency Market Risk Intelligence Platform}
\newcommand{\studentA}{Student Name 1}
\newcommand{\regA}{CIIT/SPxx-BCS-xxx/ATK}
\newcommand{\studentB}{Student Name 2}
\newcommand{\regB}{CIIT/SPxx-BCS-xxx/ATK}
\newcommand{\supervisor}{Supervisor Name}
\newcommand{\supervisordesignation}{Lecturer}
\newcommand{\hod}{Head of Department Name}
\newcommand{\hoddesignation}{Associate Professor}
\newcommand{\examiner}{External Examiner Name}
\newcommand{\examinerorg}{Examiner's Institution}
\newcommand{\session}{FALL, 2026}
\newcommand{\submissionmonth}{December, 2026}
\newcommand{\university}{COMSATS University Islamabad}
\newcommand{\campus}{Attock Campus}

% ------------------------------------------------------ captions and lists
\captionsetup{labelsep=space,font=normalsize,justification=centering}
\renewcommand{\figurename}{Figure}
\renewcommand{\tablename}{Table}
\renewcommand{\cftfigpresnum}{Fig }
\renewcommand{\cfttabpresnum}{Table }
\setlength{\cftfignumwidth}{4.5em}
\setlength{\cfttabnumwidth}{5.5em}
\renewcommand{\cftchapleader}{\cftdotfill{\cftdotsep}}
\renewcommand{\contentsname}{TABLE OF CONTENTS}
\renewcommand{\listfigurename}{LIST OF FIGURES}
\renewcommand{\listtablename}{LIST OF TABLES}
\renewcommand{\cfttoctitlefont}{\hfill\Large\bfseries}
\renewcommand{\cftaftertoctitle}{\hfill}
\renewcommand{\cftloftitlefont}{\hfill\Large\bfseries}
\renewcommand{\cftafterloftitle}{\hfill}
\renewcommand{\cftlottitlefont}{\hfill\Large\bfseries}
\renewcommand{\cftafterlottitle}{\hfill}
\setcounter{tocdepth}{3}
\setcounter{secnumdepth}{3}

% ------------------------------------------------------- heading styles
% Chapter: a page of its own with "Chapter N" and the title centred (as in the template).
\titleformat{\chapter}[display]{\centering\normalfont\bfseries\LARGE}{\chaptertitlename\ \thechapter}{18pt}{\LARGE}
\titlespacing*{\chapter}{0pt}{0.36\textheight}{0pt}  % pushes the chapter title to the middle of its page
\titleformat{\section}{\normalfont\Large\bfseries}{\thesection}{0.8em}{}
\titleformat{\subsection}{\normalfont\large\bfseries}{\thesubsection}{0.8em}{}
\titleformat{\subsubsection}{\normalfont\normalsize\bfseries}{\thesubsubsection}{0.8em}{}

% \chapter already starts a new page; the title page is then closed so the chapter text starts on the next page.
\newcommand{\fypchapter}[2]{\chapter{#1}\label{#2}\thispagestyle{empty}\clearpage}
\newcommand{\chapteroverview}{\begin{center}{\LARGE\bfseries Chapter Overview}\end{center}\vspace{4pt}}

% ---------------------------------------------------------- helpers
% \diagram{file in figures/}{caption}{label}{width}: shows the image if it exists, otherwise a placeholder box.
\newcommand{\diagram}[4]{%
  \begin{figure}[H]\centering
  \IfFileExists{figures/#1}{\includegraphics[width=#4,keepaspectratio,height=0.78\textheight]{#1}}{%
    \fbox{\parbox[c][5cm][c]{0.8\linewidth}{\centering Add figure: \texttt{figures/\detokenize{#1}}}}}
  \caption{#2}\label{#3}\end{figure}}

% Functional requirement table in the template's layout.
% \frtable{caption}{label}{id}{title}{user perspective}{system perspective}{source}{rationale}{business rule}{dependencies}{priority}
\newcommand{\frtable}[11]{%
  \begin{table}[H]\centering\caption{#1}\label{#2}
  \renewcommand{\arraystretch}{1.35}
  \begin{tabularx}{\linewidth}{|p{0.27\linewidth}|X|}\hline
  \textbf{Identifier} & #3 \\\hline
  \textbf{Title} & #4 \\\hline
  \textbf{Requirement} & \textbf{User Perspective:} #5\par\textbf{System Perspective:} #6 \\\hline
  \textbf{Source} & #7 \\\hline
  \textbf{Rationale} & #8 \\\hline
  \textbf{Business Rule (if required)} & #9 \\\hline
  \textbf{Dependencies} & #10 \\\hline
  \textbf{Priority} & #11 \\\hline
  \end{tabularx}\end{table}}

% Test case block in the template's style.
\newcommand{\testcase}[5]{%
  \textbf{#1}\par
  \textit{Testing Objective:} #2\par
  \textit{Expected Result:} #3\par
  \textit{Actual Result:} #4\par
  \textit{Status:} #5\par\medskip}

\begin{document}

% ======================================================= front matter
\pagenumbering{roman}
% ---------------------------------------------------------------- logo helper
\newcommand{\unilogo}[1]{\IfFileExists{figures/comsats_logo.png}{\includegraphics[width=#1]{comsats_logo.png}}%
  {\fbox{\parbox[c][#1][c]{#1}{\centering\small Put the university logo at\\ \texttt{figures/comsats\_logo.png}}}}}

% ================================================================ title page 1
\begin{titlepage}
\centering
\vspace*{0.4cm}
{\fontsize{26}{32}\selectfont \projecttitle\par}
\vspace{1.4cm}
\unilogo{4.2cm}\par
\vspace{1.4cm}
{\Large\itshape By\par}
\vspace{0.4cm}
{\Large \studentA\\ \regA\\ \studentB\\ \regB\par}
\vspace{1cm}
{\Large Project\par}\vspace{0.4cm}
{\Large for\par}\vspace{0.4cm}
{\Large Bachelor of Science in Computer Science\par}
\vspace{0.8cm}
{\LARGE \university\\[4pt] \campus{} - Pakistan\par}
\vspace{0.8cm}
{\Large \session\par}
\end{titlepage}
\setcounter{page}{2}

% ================================================================ title page 2
\begin{center}
\begin{minipage}{0.22\linewidth}\unilogo{2.6cm}\end{minipage}%
\begin{minipage}{0.75\linewidth}\centering{\Large\bfseries \university, \campus}\end{minipage}
\vspace{1.4cm}

{\fontsize{22}{28}\selectfont \projecttitle\par}
\vspace{1.4cm}
{\large This Project is Presented to\par}\vspace{0.8cm}
{\Large \university{} \campus\par}\vspace{0.8cm}
{\large In partial fulfillment\\ of the requirement for the degree of\par}\vspace{0.8cm}
{\LARGE BS in Computer Science\par}\vspace{0.8cm}
{\large By\par}\vspace{0.5cm}
{\large \studentA\\ \regA\\ \studentB\\ \regB\par}\vspace{1cm}
{\LARGE \session\par}
\end{center}
\clearpage

% ================================================================ submission page
\begin{center}
{\fontsize{22}{28}\selectfont \underline{CryptoSense AI: An Explainable AI-Based}\\[4pt]\underline{Cryptocurrency Market Risk Intelligence Platform}\par}
\vspace{1cm}
\begin{minipage}{0.92\linewidth}\large
This final year project is submitted to the Department of Computer Science as partial fulfillment of the requirement for the award of the degree of BS Computer Science.
\end{minipage}
\vspace{1.2cm}

{\renewcommand{\arraystretch}{2.2}
\begin{tabular}{|p{0.38\linewidth}|p{0.42\linewidth}|}\hline
\rowcolor{gray!60}\centering\textbf{Name} & \centering\arraybackslash\textbf{Registration Number}\\\hline
\centering \studentA & \centering\arraybackslash \regA \\\hline
\centering \studentB & \centering\arraybackslash \regB \\\hline
\end{tabular}}
\end{center}
\vfill
\begin{flushleft}
{\large\bfseries Supervisor\par}\vspace{0.6cm}
{\large \supervisor\\[4pt] \supervisordesignation,\\[4pt] Department of Computer Science\\[4pt] \university, \campus\par}
\vspace{0.8cm}
{\large \submissionmonth\par}
\end{flushleft}
\clearpage

% ================================================================ final approval
\begin{center}
{\LARGE\bfseries Final Approval\par}\vspace{2pt}\rule{\linewidth}{0.5pt}\vspace{0.6cm}

{\large This Project titled\par}\vspace{0.4cm}
{\fontsize{20}{26}\selectfont \projecttitle\par}\vspace{1cm}
{\large By\par}\vspace{0.8cm}
{\large\itshape \studentA\\ \regA\\ \studentB\\ \regB\par}\vspace{1cm}
{\large Has been approved for the \university, \campus.\par}
\end{center}
\vspace{0.6cm}
{\large External Examiner: \rule{0.6\linewidth}{0.4pt}\par}
\begin{center}{\large \examiner\\ \examinerorg\par}\end{center}
\vspace{0.4cm}
{\large Supervisor: \rule{0.66\linewidth}{0.4pt}\par}
\begin{center}{\large \supervisor\\ Department of Computer Science, (CUI) \campus\par}\end{center}
\vspace{0.4cm}
{\large HoD: \rule{0.75\linewidth}{0.4pt}\par}
\begin{center}{\large \hod\\ Department of Computer Science, (CUI) \campus\par}\end{center}
\clearpage

% ================================================================ declaration
\begin{center}{\LARGE\bfseries\underline{Author's Declaration}}\end{center}
\vspace{0.4cm}
We, \studentA{} and \studentB, \regA{} and \regB, hereby declare that we have produced the work presented in this project during the scheduled period of study. We also declare that we have not taken any material from any source except where it is referred to, and that the amount of plagiarism is within the acceptable range. If a violation of HEC rules on research has occurred in this thesis, we shall be liable to punishable action under the plagiarism rules of HEC.
\vspace{2.5cm}
\begin{flushright}
Student Signature\\[0.6cm]
\studentA\\ \regA\\[1cm]
Student Signature\\[0.6cm]
\studentB\\ \regB
\end{flushright}
\vspace{0.8cm}
Date: \rule{4cm}{0.4pt}
\clearpage

% ================================================================ certificate
\begin{center}{\LARGE\bfseries\underline{Certificate}}\end{center}
\vspace{0.4cm}
It is certified that \studentA, \regA{} and \studentB, \regB, have carried out all the work related to this thesis under my supervision at the Department of Computer Science, \university, \campus, and the work fulfills the requirements for the award of the degree of BS in Computer Science.
\vspace{2cm}
\begin{flushright}
Supervisor\\[1cm]\rule{4.5cm}{0.4pt}\\[0.6cm]
\supervisor\\ \supervisordesignation,\\ Department of Computer Science,\\ CUI, \campus
\end{flushright}
\vspace{0.6cm}
Head of Department:\\[1cm]\rule{4.5cm}{0.4pt}\\[0.6cm]
\hod\\ \hoddesignation,\\ Department of Computer Science,\\ CUI, \campus
\clearpage

% ================================================================ dedication
\vspace*{\fill}
\begin{center}
{\Large\itshape Dedicated to our parents,\\[8pt] who have always been our biggest supporters\\[8pt] and a constant source of inspiration\par}
\end{center}
\vspace*{\fill}
\clearpage

% ================================================================ acknowledgements
\begin{center}
{\LARGE\bfseries ACKNOWLEDGEMENTS\par}\vspace{0.4cm}
{\large\bfseries Praise be to ALLAH, the Cherisher and Lord\\ of the World, Most Gracious and Most Merciful\par}
\end{center}
\vspace{0.3cm}
We would like to express our sincere gratitude to \supervisor, our supervisor, for the invaluable guidance, continuous support and expert insights given throughout the development of our project titled ``\projecttitleinline.'' This mentorship was essential in shaping the project and bringing it to completion. We also thank the faculty of the Department of Computer Science for providing the resources and an environment in which learning and growth could take place. Finally, we give our heartfelt thanks to our families and peers, who supported us and gave us moral support throughout this work.
\vspace{1.5cm}
\begin{flushright}
\studentA\\ \regA\\[1cm]
\studentB\\ \regB
\end{flushright}
\clearpage

% ================================================================ project summary
\begin{center}{\LARGE\bfseries Project Summary}\end{center}
\vspace{0.2cm}
CryptoSense AI is a web application that helps everyday cryptocurrency users understand what the market is doing and how risky it is, in plain language. Most retail investors see prices and charts but not the context behind them: how volatile a coin is, whether leveraged traders are being forced out, what the news mood is, or how confident any prediction really is. CryptoSense brings these signals together in one place and explains them.

The system streams live prices from Binance and computes 18 features per coin from hourly candles, futures funding rates, the Fear and Greed index and Bitcoin's trend. An ensemble of XGBoost and LightGBM models classifies the next 24 hours as bullish, bearish or neutral, and a transparent seven-factor risk score shows exactly why a coin is rated low, moderate or high risk. A Risk Copilot adds model votes, similar historical market periods, observed price levels and a volatility what-if simulation.

A conversational assistant, Sage, coordinates six specialist analysts (market, derivatives, on-chain, news and sentiment, risk and explanation) and answers questions using retrieval over a knowledge base. Large language models (Anthropic Claude, with an OpenAI-compatible fallback) are used only to phrase answers; an explanation module rejects any answer that contains a number not found in the verified data. A Liquidation Intelligence module records Binance futures liquidations, computes statistics over five time windows, detects unusual activity with robust statistics, relates it to price, open interest and funding, and raises alerts by Telegram. Users can also manage a portfolio with stop-loss and take-profit limits, set alerts, and practise with a paper-trading account.

The backend is built with Python and FastAPI, the frontend with Flutter for the web, and data is stored with SQLAlchemy in SQLite or PostgreSQL. A strict no-fabricated-data policy is applied throughout: when a live input is missing, the system reports it as unavailable instead of guessing. The system is verified by 65 automated backend tests, static analysis of the frontend and live end-to-end checks. Model results are reported honestly against a majority-class baseline, showing both what the models can and cannot do.
\clearpage


\tableofcontents
\clearpage
\listoffigures
\clearpage
\listoftables
\clearpage

% ======================================================= chapters
\pagenumbering{arabic}
\fypchapter{Introduction}{ch:intro}

\chapteroverview
This chapter introduces CryptoSense AI, an explainable, AI-based cryptocurrency market risk intelligence platform. It explains the motivation behind the system, its significance, the problem it addresses, its scope and objectives, its limitations, the project deliverables, its relevance to the courses of the BS Computer Science programme, and the research objectives and questions that guided the work. It closes with the contributions of the project.

\section{Introduction}
Cryptocurrencies such as Bitcoin and Ethereum trade twenty-four hours a day, and their prices often move several percent within hours. Millions of retail users now hold crypto assets, but most of them see only a price and a chart. The factors that actually describe the risk of a position, such as volatility, leverage in the futures market, forced liquidations, large transfers by ``whales'', news sentiment and the behaviour of Bitcoin itself, are spread across many specialist websites, are presented in technical language, and are rarely explained.

CryptoSense AI is a web application that brings these signals together and explains them in plain words. It streams live market data from Binance, computes technical and market-context features, and uses an ensemble of gradient-boosted machine learning models (XGBoost and LightGBM) to classify the next 24 hours for each tracked coin. A transparent seven-factor risk score shows exactly why a coin is rated low, moderate or high risk, and a Risk Copilot adds model agreement, similar historical periods, observed price levels and a volatility what-if simulation.

A conversational assistant named Sage coordinates six specialist analysts and answers user questions with retrieval-augmented generation. Large language models are used only to phrase verified results: the system checks every number in an AI-written explanation against the computed data and rejects the answer if a number cannot be traced. A Liquidation Intelligence module records forced closures of leveraged Binance futures positions, detects unusual activity with robust statistics, relates it to price, open interest and funding, and sends alerts. Users can also track a portfolio with stop-loss and take-profit limits, create alerts, read news with sentiment, monitor whale activity and practise trading with pretend money.

A central design principle of the project is honesty about data and models. When a live input is unavailable, CryptoSense reports it as unavailable instead of substituting made-up values, and model results are always reported against a simple baseline so that users and examiners can see what the models can and cannot do.

\section{Significance of the Proposed System}
\textbf{Vision Statement:} For retail cryptocurrency users, students and beginner investors who need to understand market conditions and risk before making decisions, CryptoSense AI is a market risk intelligence platform that combines live market data, machine learning and explainable analytics in one place. Unlike price trackers that show numbers without context, or professional analytics platforms that are expensive and written for experts, CryptoSense explains what is happening and why, states the uncertainty of every model output, and never presents missing or simulated data as live evidence.

The system is significant for three reasons. First, it lowers the barrier to understanding market risk: concepts such as funding rates, liquidations and volatility are explained where they appear. Second, it promotes responsible behaviour: the platform does not execute trades or tell users to buy or sell, and it encourages stop-losses and position sizing. Third, it demonstrates a trustworthy way to use large language models in a financial setting, where every figure in an AI-written explanation is verified against computed data.

\section{Problem Statement}
Retail cryptocurrency users make decisions with incomplete and fragmented information. Prices are easy to find, but the signals that describe risk are spread across different platforms, require technical knowledge and are often behind paid subscriptions. Users rarely know how volatile an asset currently is, whether leveraged traders are being liquidated, how strongly futures traders lean in one direction, or what the news mood is. Many applications that do show predictions present them as confident signals without stating their accuracy, and some fill gaps with simulated data that looks real. Chat assistants based on large language models can make the situation worse by inventing prices or statistics.

There is therefore a need for a single system that (i) collects live market, derivatives, liquidation, sentiment and on-chain data, (ii) turns it into transparent, explainable risk information, (iii) applies machine learning with honest evaluation, and (iv) explains everything in plain language without inventing numbers.

\section{Scope and Objectives}
CryptoSense AI is a web application with a Python backend that analyses six major cryptocurrencies (BTC, ETH, BNB, SOL, XRP and ADA) traded against USDT on Binance. Its scope includes live market monitoring, machine learning market-state signals, an explainable risk score and Risk Copilot, a conversational assistant with six specialist analysts, liquidation intelligence for Binance USD-M futures, news sentiment, whale and smart-money tracking (when the optional data keys are configured), new listing discovery, portfolio tracking with profit and loss limits, configurable alerts with Telegram delivery, paper trading and an educational Academy. Order execution, custody of funds and financial advice are outside the scope of the project.

\textbf{Project Objectives}
\begin{enumerate}[label=PO\arabic*]
\item To build a live market monitoring platform that streams real prices and candles from Binance and never displays fabricated market data.
\item To develop and honestly evaluate a machine learning ensemble (XGBoost and LightGBM) that classifies the next 24 hours of each tracked coin, using a leak-free chronological evaluation and a majority-class baseline.
\item To provide an explainable risk score and a Risk Copilot that show the factors behind every score, model agreement, similar historical periods and what-if scenarios.
\item To build a conversational assistant that coordinates six specialist analysts and uses language models only to phrase verified results.
\item To build a liquidation intelligence module that records observed Binance liquidations, computes multi-window statistics, detects unusual activity and raises alerts.
\item To give users tools to manage risk: portfolio tracking, stop-loss and take-profit limits, alerts with Telegram delivery and a paper-trading account.
\end{enumerate}

\section{System Limitations/Constraints}
\textbf{[L1]} The machine learning models are trained on public market data only. Crypto markets are noisy and largely efficient, so the models provide weak but measurable signals; on recent unseen data their accuracy is below that of always predicting ``neutral'' (see Chapter~\ref{ch:testing}). Their outputs are presented as model scores, not guarantees.

\textbf{[L2]} The ensemble probabilities are not yet calibrated, so the bearish score shown in the risk factors is labelled as a raw model score and not as a probability.

\textbf{[L3]} Binance's liquidation stream sends at most one liquidation per symbol per second. Liquidation figures are therefore observed samples of activity, not exchange-wide totals, and history only exists from the moment the backend starts recording.

\textbf{[L4]} Whale transfers, congressional and insider trades of crypto-linked stocks, LLM-written answers, Jev news judgement, CryptoPanic news and Telegram delivery require optional API keys. Without them these features report that they are unavailable.

\textbf{[L5]} Futures funding history on Binance begins in 2019--2020 and the Fear and Greed index in 2018, so models can only be trained on about six to seven years of complete, real data even when ten years are requested.

\textbf{[L6]} The seven-factor risk score uses documented but hand-chosen weights and level boundaries; these boundaries have not yet been validated against realised outcomes.

\textbf{[L7]} The system depends on the availability of third-party APIs. When they are unreachable, the affected features show ``unavailable'' instead of data.

\textbf{[L8]} CryptoSense is an analysis and education tool. It does not execute orders and does not provide financial advice.

\section{Project Deliverables}
The major deliverables of the CryptoSense AI project are:
\begin{itemize}
\item A complete web application (Flutter) with eleven screens: Overview, Portfolio, Alerts, Profit and Loss Limits, News, Whale Watch, Liquidations, New Listings, Academy, Risk Copilot chat and Paper Trading.
\item A backend REST API (FastAPI) with 47 endpoints, background workers and a relational database.
\item A trained XGBoost and LightGBM ensemble with a reproducible training pipeline and a metrics report.
\item A Liquidation Intelligence module with statistics, anomaly detection, market context, explanations and alerts.
\item An automated test suite of 65 backend tests.
\item Technical documentation, design diagrams, a project report and a presentation.
\end{itemize}

\section{Relevance to Course Modules}
CryptoSense AI applies knowledge from many courses of the BS Computer Science programme. \textbf{Programming Fundamentals} and \textbf{Object-Oriented Programming} were used to build modular Python services and Dart widgets. \textbf{Data Structures and Algorithms} and \textbf{Design and Analysis of Algorithms} guided the rolling-window statistics, de-duplication with bounded queues and efficient nearest-neighbour search for historical market twins. \textbf{Software Engineering} practices were followed for requirements, design, testing and documentation.

\textbf{Database Systems} concepts were applied in the relational schema, indexes for time-series queries and SQL aggregation. \textbf{Computer Networks} knowledge was needed for REST, WebSocket streaming, reconnection and timeouts. \textbf{Information Security} principles were applied through password hashing with bcrypt, JSON Web Tokens and keeping secrets out of source control. \textbf{Artificial Intelligence}, \textbf{Machine Learning}, and \textbf{Probability and Statistics} were central to feature engineering, gradient-boosted classifiers, chronological evaluation, calibration, robust z-scores and retrieval-augmented generation. \textbf{Human--Computer Interaction} shaped the plain-language, beginner-friendly interface, and \textbf{Parallel and Distributed Computing} concepts were applied in the asynchronous background workers.

\section{Research Objective}
\begin{itemize}
\item To combine live market, derivatives, liquidation, sentiment and on-chain signals into a single explainable view of cryptocurrency market risk.
\item To measure how well gradient-boosted models trained on public data can classify short-horizon market states, using a leak-free chronological evaluation.
\item To study how the amount of training history (two, five and ten years requested) affects out-of-sample performance on the same unseen period.
\item To design a safe way of using large language models for financial explanations, in which every number is verified against computed data.
\item To detect unusual liquidation activity with robust statistical baselines rather than fixed thresholds.
\end{itemize}

\section{Research Questions}
\begin{enumerate}[label=RQ\arabic*.]
\item How accurately can XGBoost and LightGBM classify the next 24 hours of major cryptocurrencies from public market data, compared with a majority-class baseline?
\item Does training on a longer history improve performance on the same recent, unseen period?
\item Is short-horizon downside risk (the size of moves) more predictable than direction?
\item How can a large language model be used to explain market analytics without inventing prices, amounts or probabilities?
\item How can unusual liquidation activity be identified reliably from a sampled stream?
\end{enumerate}

\section{Contribution}
This project contributes a working, end-to-end market risk intelligence platform that combines machine learning, transparent risk scoring, liquidation analytics and a multi-specialist conversational assistant. Its main contributions are: (1) an explainable risk score whose factors are always visible; (2) a grounded explanation layer that verifies every number written by a language model; (3) a liquidation intelligence module that clearly separates observed stream data from market-wide totals and uses robust statistics for spike detection; (4) a strict no-fabricated-data policy across the backend and frontend; and (5) an honest evaluation of the machine learning models, including a comparison of training-history lengths and an experiment on downside-risk targets.

\section{Chapter Summary}
This chapter introduced CryptoSense AI and explained the problem of fragmented, hard-to-understand market risk information faced by retail cryptocurrency users. It presented the vision, objectives and scope of the system, its limitations, deliverables and relevance to the degree programme, and the research objectives and questions that guided the project. The next chapter reviews existing platforms and research related to the proposed system.

\fypchapter{Literature Review}{ch:literature}

\chapteroverview
This chapter reviews existing cryptocurrency platforms and research related to market analysis, risk, liquidations, on-chain activity, sentiment and machine learning prediction. It identifies the limitations of existing solutions: fragmented information, technical presentation, paid access, predictions without honest evaluation, and AI assistants that may invent figures. It then shows how CryptoSense AI addresses these gaps.

\section{Market Data Aggregators (CoinMarketCap and CoinGecko)}
Market data aggregators such as CoinMarketCap and CoinGecko list prices, market capitalisation, trading volume and price history for thousands of assets. They are the most common starting point for retail users. However, they focus on reporting numbers rather than interpreting them: they do not explain current risk, do not combine derivatives, liquidation and sentiment signals into one assessment, and do not provide machine learning analysis with stated accuracy. CryptoSense uses live data as a starting point and adds explainable risk analysis on top of it.

\section{Charting Platforms (TradingView)}
TradingView provides professional charting with a large library of technical indicators and community scripts. It is powerful but assumes that the user already knows how to read indicators such as RSI, MACD and Bollinger Bands, and many advanced features require a subscription. It does not explain in plain language what the indicators mean for risk, and community signals are not evaluated against a baseline. CryptoSense computes the indicators automatically, explains them in plain words (through Sage and the Academy) and evaluates its model signals honestly.

\section{Derivatives and Liquidation Analytics (CoinGlass)}
Derivatives analytics platforms such as CoinGlass publish liquidation totals, open interest, funding rates and long/short ratios across exchanges. These data are valuable for understanding leverage in the market, but they are presented as dense dashboards for experienced traders, and some datasets require paid plans. They also rarely explain what a liquidation spike means or how it relates to price and open interest. CryptoSense builds its own liquidation intelligence from Binance's public stream, labels figures clearly as observed stream data, detects unusual activity with robust statistics and explains possible scenarios with the evidence that supports or contradicts each one.

\section{On-Chain and Whale Tracking (Whale Alert and Glassnode)}
Whale Alert reports large blockchain transfers in near real time and labels many exchange wallets, while on-chain analytics platforms such as Glassnode provide metrics derived from blockchain data. These services are useful but are either limited to raw transfer notifications or aimed at professional analysts, and full access usually requires a paid plan. A common mistake in retail interpretation is to treat every exchange inflow as selling. CryptoSense integrates Whale Alert when a key is configured, classifies exchange inflow and outflow, and explains transfers as observations rather than proof of intent.

\section{Sentiment Platforms (LunarCrush and Santiment)}
Social and news sentiment platforms track mentions, sentiment and attention across social media and news sources. They show that sentiment is an important context for crypto prices, but their scores are often proprietary and difficult to interpret. CryptoSense scores each headline with VADER, a transparent lexicon-based sentiment model \cite{hutto2014vader} extended with a crypto vocabulary, optionally compares it with the Jev decision model, and combines news sentiment with the Fear and Greed index.

\section{Machine Learning for Cryptocurrency Prediction}
Many studies apply machine learning to cryptocurrency price prediction. McNally et al. \cite{mcnally2018} compared recurrent neural networks and LSTM models for Bitcoin price direction and reported modest accuracy. Chen et al. \cite{chen2020bitcoin} showed that the best model depends on data frequency and that simpler statistical models can perform as well as complex ones on daily data. Gradient-boosted decision trees such as XGBoost \cite{chen2016xgboost} and LightGBM \cite{ke2017lightgbm} are strong on tabular indicator data and are widely used in practice.

A recurring problem in this literature is evaluation. Random train/test splits allow information from the future to leak into training, which produces unrealistically high accuracy \cite{lopezdeprado2018}. The efficient market hypothesis \cite{fama1970} suggests that publicly available information is quickly reflected in prices, so modest accuracy is expected. Probabilistic forecasts should also be judged with proper scoring rules such as the Brier score \cite{brier1950} and calibrated before being presented as probabilities \cite{niculescu2005}. CryptoSense follows these lessons: it uses a chronological split with gaps, compares every model with a majority-class baseline, and labels uncalibrated outputs as raw scores.

\section{Grounded Conversational Assistants}
Large language models can explain complex topics in plain language, but they can also produce fluent statements that are not supported by data. Retrieval-augmented generation \cite{lewis2020rag} reduces this risk by giving the model relevant source passages and asking it to cite them. In a financial application, an invented price or probability could mislead a user, so retrieval alone is not enough. CryptoSense combines retrieval with a verification step: every number in an AI-written explanation must match a number in the computed facts, otherwise the answer is replaced by a deterministic, rules-based explanation.

\section{Comparison with Existing Solutions}
Table~\ref{tab:related} summarises the weaknesses of existing solutions and how CryptoSense AI addresses them.

{\renewcommand{\arraystretch}{1.3}
\begin{longtable}{|p{0.2\linewidth}|p{0.34\linewidth}|p{0.36\linewidth}|}
\caption{Related System Analysis with Proposed CryptoSense AI Solution}\label{tab:related}\\\hline
\textbf{Application} & \textbf{Weakness} & \textbf{Proposed Project Solution}\\\hline\endfirsthead
\hline\textbf{Application} & \textbf{Weakness} & \textbf{Proposed Project Solution}\\\hline\endhead
CoinMarketCap / CoinGecko & Report prices and volumes but do not interpret risk or combine signals. & Combines live prices with derivatives, liquidation, sentiment and ML signals into an explained risk view.\\\hline
TradingView & Assumes indicator knowledge; advanced features are paid; community signals are not evaluated. & Computes indicators automatically, explains them in plain words and evaluates model signals against a baseline.\\\hline
CoinGlass & Dense expert dashboards; some data paid; little explanation of what spikes mean. & Observed liquidation statistics, robust spike detection and scenario explanations with supporting and contradicting evidence.\\\hline
Whale Alert / Glassnode & Raw notifications or expert metrics; paid access; inflows often misread as selling. & Classifies exchange flows and explains them as observations, not proof of intent.\\\hline
LunarCrush / Santiment & Proprietary sentiment scores that are hard to interpret. & Transparent VADER sentiment with a crypto lexicon, optional Jev comparison and the Fear and Greed index.\\\hline
General AI chatbots & May invent prices, statistics or probabilities. & Six specialist analysts, retrieval with citations, and a verification step that rejects unverified numbers.\\\hline
\end{longtable}}

\section{Chapter Summary}
This chapter reviewed market data aggregators, charting platforms, derivatives and liquidation analytics, on-chain and whale trackers, sentiment platforms, machine learning research and grounded conversational assistants. Existing solutions are fragmented, technical, often paid, and rarely evaluate predictions honestly. CryptoSense AI addresses these gaps with a single explainable platform, honest model evaluation and a verified AI explanation layer. The next chapter specifies the requirements of the system.

\fypchapter{Requirement Analysis}{ch:requirements}

\chapteroverview
This chapter specifies the functional and non-functional requirements of CryptoSense AI. It describes the user classes, the techniques used to identify requirements, twelve functional requirements, the quality attributes the system must satisfy, and the user, software, hardware and communication interfaces.

\section{User Classes and Characteristics}
\begin{itemize}
\item \textbf{Beginner investors and learners:} people who are new to cryptocurrency and want to understand prices, risk and basic indicators. They need plain language, guidance (the Sage tour, lightbulb explanations and the Academy) and protection from over-confident signals.
\item \textbf{Active retail traders:} users who follow the market daily and want consolidated risk information, liquidation and derivatives context, alerts and portfolio limits. They value speed, accuracy and Telegram notifications.
\item \textbf{Students and researchers:} users who study market behaviour and machine learning. They need transparent methods, reproducible training and honest metrics (available through the model-information endpoint and this documentation).
\item \textbf{Operator / deployer:} the person who runs the backend (locally or on Railway), configures optional API keys and Telegram delivery, and trains the models.
\end{itemize}

\section{Requirements Identification Technique}
\begin{itemize}
\item \textbf{Comparative study:} existing platforms (Chapter~\ref{ch:literature}) were analysed to identify missing features, such as plain-language risk explanations and honest model evaluation.
\item \textbf{Use case modelling:} the main interactions of guests, registered users and external systems were modelled as use cases (Figure~\ref{fig:usecase}).
\item \textbf{Prototyping and iteration:} the interface and analytics were built incrementally, with each feature reviewed in the running application before the next one was added.
\item \textbf{Specification documents:} detailed written specifications were prepared for the liquidation tracker, the risk intelligence specialist and the orchestrator. They defined data-integrity rules such as ``never use fabricated data'' and ``label observed data clearly''.
\end{itemize}

\section{Functional Requirements}

\subsection{Functional Requirement 1 -- Live Market Monitoring}
\frtable{Description of FR-1 - Live market monitoring}{tab:fr1}{FR-1}{Live Market Monitoring}
{The user shall see live prices, 24-hour change, high, low and volume for the tracked coins, and an interactive candlestick chart on 15-minute, 1-hour, 4-hour and 1-day intervals.}
{The system shall load a real snapshot from Binance REST, stream live tickers over WebSocket, and show ``unavailable'' with automatic retry if Binance cannot be reached.}
{Comparative study of market data platforms; data-integrity specification.}
{Live, trustworthy prices are the basis of every other feature.}
{No simulated or hardcoded prices may be shown in place of real data.}
{Binance public API.}
{High}

\subsection{Functional Requirement 2 -- AI Market Signals}
\frtable{Description of FR-2 - AI market signals}{tab:fr2}{FR-2}{AI Market Signals}
{The user shall see, for each tracked coin, whether the next 24 hours are classified as bullish, bearish or neutral, with the model score, the individual model votes and the reasons.}
{The system shall build 18 features from real hourly candles, futures funding, the Fear and Greed index and Bitcoin's trend, and run the XGBoost and LightGBM ensemble. If any input is missing or no trained model exists, it shall report the signal as unavailable and list the missing inputs.}
{Machine learning research review; project objectives PO2.}
{Gives users a second opinion that is transparent about its inputs and limitations.}
{Directional calls with a top score below 0.55 are shown as neutral.}
{FR-1, trained model files.}
{High}

\subsection{Functional Requirement 3 -- Explainable Risk Score and Risk Copilot}
\frtable{Description of FR-3 - Explainable risk score and Risk Copilot}{tab:fr3}{FR-3}{Explainable Risk Score and Risk Copilot}
{The user shall see a 0--100\% risk score and level for each coin, the seven factors behind it, model agreement, similar historical periods, observed key levels and a volatility what-if simulation.}
{The system shall compute the score from seven weighted factors (price swings 30, trend 15, AI model view 15, momentum stretch 10, crowded trades 10, market mood 10, Bitcoin trend 10), store a history at most every 30 minutes, find historical twins with nearest-neighbour search, and recompute the score for a changed volatility.}
{Problem statement; user need for understandable risk.}
{Users can see why a coin is risky instead of receiving an unexplained number.}
{Levels: below 35 low, 35--59 moderate, 60 or above high. Historical twins are descriptive, not forecasts.}
{FR-2.}
{High}

\subsection{Functional Requirement 4 -- Conversational Risk Copilot (Sage)}
\frtable{Description of FR-4 - Conversational Risk Copilot}{tab:fr4}{FR-4}{Conversational Risk Copilot (Sage)}
{The user shall ask questions in English, Urdu or Roman Urdu and receive a short, beginner-friendly answer with sources and the status of six specialist analysts.}
{The system shall gather live context, run six specialists (market, derivatives, on-chain, sentiment and news, risk, explanation), retrieve relevant knowledge with TF-IDF search, and generate the answer with Claude, then an OpenAI-compatible model, then rules, in that order of availability.}
{Problem statement; grounded assistant research.}
{Lets beginners ask about the market in their own words.}
{The assistant never tells the user to buy or sell and treats unavailable data literally.}
{FR-1, FR-2, FR-3, FR-6, FR-7.}
{High}

\subsection{Functional Requirement 5 -- Liquidation Intelligence}
\frtable{Description of FR-5 - Liquidation intelligence}{tab:fr5}{FR-5}{Liquidation Intelligence}
{The user shall view observed liquidations for 5-minute, 15-minute, 1-hour, 4-hour and 24-hour windows, filtered by coin, with totals, long/short split and ratio, the largest event, the most affected coin, charts, spike detection, market context and an AI explanation.}
{The system shall consume the Binance forceOrder stream, validate and de-duplicate events, classify forced SELL orders as long liquidations and forced BUY orders as short liquidations, store them for seven days, and compute statistics, robust-z anomalies and context from price, open interest and funding.}
{Liquidation tracker specification.}
{Liquidations reveal forced deleveraging that often accompanies sharp moves.}
{All figures are labelled as observed stream data; a spike needs both a robust z-score of at least 4 and at least \$250{,}000 of activity (configurable).}
{Binance USD-M futures API; FR-9.}
{High}

\subsection{Functional Requirement 6 -- News and Sentiment}
\frtable{Description of FR-6 - News and sentiment}{tab:fr6}{FR-6}{News and Sentiment}
{The user shall read recent crypto headlines tagged bullish, bearish or neutral, with the coins mentioned and Sage's take, and filter them by sentiment.}
{The system shall collect headlines from RSS feeds (or CryptoPanic when configured) every 15 minutes, score them with VADER and a crypto lexicon (or Jev when configured), and add them to the knowledge base.}
{Comparative study of sentiment platforms.}
{News often moves prices; sentiment adds context to the risk view.}
{If no feed is reachable, news is reported as unavailable; no sample headlines are shown.}
{Internet access to news sources.}
{Medium}

\subsection{Functional Requirement 7 -- Whale and Smart-Money Tracking}
\frtable{Description of FR-7 - Whale and smart-money tracking}{tab:fr7}{FR-7}{Whale and Smart-Money Tracking}
{The user shall see large on-chain transfers with exchange inflow and outflow totals, and trades by members of Congress and company insiders in crypto-linked stocks.}
{The system shall query Whale Alert and Quiver Quantitative when their keys are configured and classify each transfer as exchange inflow, outflow or wallet-to-wallet.}
{Comparative study of on-chain platforms.}
{Large holders can influence supply on exchanges.}
{Transfers are presented as observations, not as proof of buying or selling. Without keys the feeds are unavailable.}
{Whale Alert and Quiver API keys (optional).}
{Medium}

\subsection{Functional Requirement 8 -- Portfolio and Profit/Loss Limits}
\frtable{Description of FR-8 - Portfolio and profit/loss limits}{tab:fr8}{FR-8}{Portfolio and Profit/Loss Limits}
{The user shall add holdings with amount and average buy price, see their value and profit or loss at live prices, and set a stop-loss and take-profit percentage per coin.}
{The system shall value holdings at live prices, flag holdings whose live price is unavailable, and notify the user when a limit is reached.}
{User need for risk management.}
{Encourages deciding exit points before emotions do.}
{A holding without a live price is clearly flagged instead of silently valued at its buy price.}
{FR-1, FR-9, FR-11.}
{High}

\subsection{Functional Requirement 9 -- Alerts and Notifications}
\frtable{Description of FR-9 - Alerts and notifications}{tab:fr9}{FR-9}{Alerts and Notifications}
{The user shall create alerts when a price rises above or falls below a value, moves by a percentage in 24 hours, when risk reaches a threshold, or for liquidation events (large liquidation, spike, long spike, short spike, imbalance shift), and receive in-app and Telegram notifications.}
{The system shall check active alerts every 30 seconds; price and risk alerts fire once, while liquidation alerts re-arm after a configurable cooldown (30 minutes by default).}
{Liquidation tracker specification; user need.}
{Users do not have to watch the market all the time.}
{Liquidation alerts require login because they are evaluated by the backend.}
{FR-1, FR-3, FR-5, FR-11; Telegram Bot API (optional).}
{High}

\subsection{Functional Requirement 10 -- Paper Trading}
\frtable{Description of FR-10 - Paper trading}{tab:fr10}{FR-10}{Paper Trading}
{The user shall practise buying and selling with 10{,}000 pretend dollars at real live prices and review the trade history and statistics.}
{The system shall apply a 0.1\% fee per trade, record each trade with the risk level and trend at the time, and pause trading when live prices are unavailable.}
{User need for safe practice.}
{Lets beginners learn without risking money.}
{Trades are only executed at live prices.}
{FR-1, FR-3, FR-11.}
{Medium}

\subsection{Functional Requirement 11 -- User Accounts and Authentication}
\frtable{Description of FR-11 - User accounts and authentication}{tab:fr11}{FR-11}{User Accounts and Authentication}
{The user shall register, log in and stay logged in, so that the portfolio, alerts and paper account are saved.}
{The system shall store passwords as bcrypt hashes, issue JSON Web Tokens valid for seven days, and protect personal endpoints with bearer authentication.}
{Security requirements.}
{Personal data must be private and persistent.}
{Guests can use all public analytics without an account.}
{Database.}
{High}

\subsection{Functional Requirement 12 -- New Listings and Academy}
\frtable{Description of FR-12 - New listings and Academy}{tab:fr12}{FR-12}{New Listings and Academy}
{The user shall see coins listed on Binance in the last 60 days with their change since listing and a risk rating, and learn market concepts step by step in the Academy.}
{The system shall refresh listings every six hours from Binance and rate each listing by volatility and volume; Academy content is also loaded into the assistant's knowledge base.}
{User need for education and discovery.}
{New coins are especially risky; education reduces mistakes.}
{Listing risk ratings are descriptive and not recommendations.}
{Binance public API.}
{Medium}

\section{Non-Functional Requirements}
\subsection{Reliability and Data Integrity}
REL-1: The system shall never display fabricated, simulated or hardcoded market values as live data. Missing inputs shall be reported as unavailable.\\
REL-2: The liquidation stream shall reconnect automatically after an error or 300 seconds without messages, and record the disconnected period.\\
REL-3: Stored liquidation history shall survive backend restarts, and statistics shall state whether a time window was fully observed.\\
REL-4: Machine learning evaluation shall use a chronological split with gaps, and training shall use only hours with real (not filled-in) context data.

\subsection{Usability}
USE-1: Every analytical panel shall include a plain-language explanation or a lightbulb help button.\\
USE-2: A first-time user shall be guided by a short Sage tour.\\
USE-3: The assistant shall answer in the user's language (English, Urdu or Roman Urdu).

\subsection{Performance}
PER-1: Live prices shall refresh at most about once per second in the interface (WebSocket updates are throttled).\\
PER-2: Signals shall be cached for five minutes and prices for five seconds to avoid repeated external calls.\\
PER-3: Liquidation statistics shall be computed with SQL aggregation and indexed queries so that 24-hour windows remain responsive.

\subsection{Security}
SEC-1: Passwords shall be stored only as bcrypt hashes.\\
SEC-2: Personal endpoints shall require a valid JSON Web Token.\\
SEC-3: API keys and tokens shall be read from environment variables and never committed to source control or shown in the interface.\\
SEC-4: Text inside market data, headlines and specialist reports shall be treated as data, not as instructions, by the language model.

\section{External Interface Requirements}
\subsection{User Interfaces Requirements}
\begin{itemize}
\item UI-1: The interface shall use a consistent dark theme with a left navigation sidebar grouped into Markets, Your money, Discover, Learn and Practice.
\item UI-2: The top bar shall show the page title, live BTC and ETH prices, a Live/Offline status badge, the time and a sign-in button.
\item UI-3: Charts shall show tooltips on hover and support time-window and coin filters.
\item UI-4: Error and unavailable states shall use simple language, for example ``Live market data unavailable''.
\item UI-5: The layout shall adapt to narrow screens by moving the sidebar into a drawer.
\end{itemize}

\subsection{Software Interfaces}
SI-1: Backend: Python 3.11, FastAPI, Uvicorn, SQLAlchemy 2, Pydantic 2.\\
SI-2: Machine learning: XGBoost, LightGBM, scikit-learn, pandas, NumPy.\\
SI-3: Database: SQLite for local use and PostgreSQL (psycopg) for deployment.\\
SI-4: Frontend: Flutter (Dart) for the web, with the http, web\_socket\_channel, shared\_preferences, fl\_chart, intl and google\_fonts packages.\\
SI-5: Language models: Anthropic Python SDK (Claude) and any OpenAI-compatible chat-completions API (Groq by default).

\subsection{Hardware Interfaces}
HI-1: The client requires a modern web browser on a desktop, laptop or mobile device.\\
HI-2: The backend requires a computer or cloud instance with at least 4 GB RAM; model training benefits from 8 GB RAM.

\subsection{Communications Interfaces}
CI-1: The client and backend communicate over HTTPS using JSON.\\
CI-2: Live prices and liquidations are received over secure WebSockets (WSS) from Binance.\\
CI-3: Alerts are delivered to Telegram through the Telegram Bot API over HTTPS when configured.

\section{Chapter Summary}
This chapter identified the user classes of CryptoSense AI and specified twelve functional requirements covering market monitoring, AI signals, explainable risk, the conversational copilot, liquidation intelligence, news, whale tracking, portfolio limits, alerts, paper trading, authentication and listings. It also defined reliability, data-integrity, usability, performance and security requirements, and the external interfaces of the system. The next chapter presents the design and architecture that satisfy these requirements.

\fypchapter{Design and Architecture}{ch:design}

\chapteroverview
This chapter describes the design and architecture of CryptoSense AI. It presents the modules of the system and their features, the layered client-server architecture, the design models (activity, use case, class, sequence and state diagrams), the data and machine learning design, the human interface, the tools and external APIs used, the REST API, and the deployment of the system.

The system consists of a Flutter web client and a FastAPI backend. The backend is organised into the modules below; each one is listed with its main features (FE).

\subsection*{Module 1: Live Market Module}
FE1 Load a real snapshot of 24-hour tickers and recent hourly closes from Binance REST.\\
FE2 Stream live tickers over WebSocket and throttle interface updates to about once per second.\\
FE3 Show interactive candlestick charts on 15m, 1h, 4h and 1d intervals.\\
FE4 Report ``unavailable'' and retry automatically when Binance cannot be reached.

\subsection*{Module 2: AI Signal Module}
FE1 Build 18 features from candles, funding, the Fear and Greed index and Bitcoin's trend.\\
FE2 Predict bullish, bearish or neutral for the next 24 hours with the XGBoost and LightGBM ensemble.\\
FE3 Show model votes and plain-language reasons; downgrade low-confidence directional calls to neutral.\\
FE4 Report the signal as unavailable, with the missing inputs, when any input is not real.

\subsection*{Module 3: Risk Engine and Risk Copilot}
FE1 Compute a 0--100\% risk score from seven transparent factors and a level.\\
FE2 Store risk history and compute value-weighted portfolio risk.\\
FE3 Find historical market twins and observed key levels (previous day, Asia and London sessions, prior week).\\
FE4 Simulate the score under a different volatility (what-if).

\subsection*{Module 4: Sage Conversational Copilot}
FE1 Run six specialist analysts and return their status and summaries.\\
FE2 Retrieve relevant knowledge (lessons, news, signals) with TF-IDF search.\\
FE3 Generate a cited answer with Claude, an OpenAI-compatible model or rules.\\
FE4 Support English, Urdu and Roman Urdu.

\subsection*{Module 5: Liquidation Intelligence Module}
FE1 Consume, validate, classify and de-duplicate Binance forceOrder events.\\
FE2 Store events for seven days and compute 5m--24h statistics and timelines.\\
FE3 Detect spikes, bursts and imbalance shifts with robust z-scores.\\
FE4 Relate activity to price, open interest, funding and volatility and explain it with verified numbers.

\subsection*{Module 6: News, Whale and Smart-Money Module}
FE1 Collect headlines and score sentiment (VADER with a crypto lexicon, optional Jev).\\
FE2 Track large transfers (Whale Alert) and classify exchange flows.\\
FE3 Show Congress and insider trades in crypto-linked stocks (Quiver Quantitative).\\
FE4 Scan Binance for new listings and rate their risk.

\subsection*{Module 7: Portfolio, Limits and Paper Trading Module}
FE1 Track holdings and profit or loss at live prices.\\
FE2 Set stop-loss and take-profit limits per coin.\\
FE3 Practise trading with 10{,}000 pretend dollars and a 0.1\% fee.

\subsection*{Module 8: Alerts and Notifications Module}
FE1 Create price, move, risk and liquidation alerts.\\
FE2 Check alerts every 30 seconds and store notifications.\\
FE3 Deliver notifications to Telegram, with a cooldown for recurring liquidation alerts.

\subsection*{Module 9: Authentication Module}
FE1 Register and log in with email and password (bcrypt).\\
FE2 Issue and verify JSON Web Tokens.

\section{Architectural Design}
CryptoSense AI follows a layered client-server architecture (Figure~\ref{fig:architecture}). The \textbf{presentation layer} is a Flutter \cite{flutter} web application that renders the eleven screens and also connects directly to Binance for live prices and candles. The \textbf{API layer} is a FastAPI \cite{fastapi} application that exposes REST endpoints grouped into routers. The \textbf{service and analytics layer} contains the feature engineering, the machine learning predictor, the risk engine, the specialist agents, the assistant with its knowledge base, the grounded explainer and the liquidation engine. \textbf{Background workers} run as asynchronous tasks inside the same process: the alert watcher, news ingestion, listings refresh, the liquidation stream and the liquidation store. The \textbf{data layer} holds the relational database, the trained model files and the knowledge store. External services are reached over HTTPS and WebSockets.

\subsection{Architecture Style}
The system combines three styles. The \textbf{client-server} style separates the user interface from the analytics and data. The \textbf{layered} style keeps routers thin: they validate input and call services, which in turn use the data layer, so analytics can be tested without the web layer. The \textbf{event-driven} style is used for liquidations: events arrive continuously from a WebSocket, are buffered in memory, written to the database in batches every three seconds, and evaluated by the alert watcher every thirty seconds. Every module returns an explicit ``unavailable'' state when its inputs are missing, so failures are visible instead of silent.

\diagram{diagrams/01_system_architecture.png}{System Architecture of CryptoSense AI}{fig:architecture}{\linewidth}

\section{Design Models}
\subsection{Activity Diagram}
Figure~\ref{fig:activity} shows the main flow of the application. When the app opens it loads a real Binance snapshot and connects to the live ticker stream; if Binance is unreachable it shows that live data is unavailable and retries after fifteen seconds. It then checks the backend. If the backend is online, signals, risk, news, whales, listings and liquidations are loaded; otherwise the backend features are shown as unavailable and the portfolio works in guest memory. From the dashboard the user can explore the Risk Copilot, ask Sage, analyse liquidations or read news in parallel, while managing money (portfolio, limits, alerts and paper trades) requires logging in. In the background the watcher checks alerts every thirty seconds.

\diagram{diagrams/03_activity_diagram.png}{CryptoSense AI - Activity Diagram}{fig:activity}{0.8\linewidth}

\subsection{Use Case Diagram}
Figure~\ref{fig:usecase} shows the use cases. A guest user can view the market, AI signals and risk, explore the Risk Copilot, chat with Sage, read news, view whale and liquidation activity, view new listings, learn in the Academy and register. A registered user inherits these use cases and can also manage a portfolio, set limits, create alerts, receive notifications and practise paper trading. Binance, the language model provider and Telegram are external actors.

\diagram{diagrams/02_use_case_diagram.png}{CryptoSense AI - Use Case Diagram}{fig:usecase}{0.78\linewidth}

\subsection{Class Diagram}
Figure~\ref{fig:class} shows the main classes. The persistent entities are \texttt{User}, \texttt{Holding}, \texttt{Alert}, \texttt{Notification}, the paper-trading classes, \texttt{RiskPrediction}, \texttt{Listing} and \texttt{LiquidationEvent}. The analytics are provided by \texttt{Predictor}, \texttt{SignalService} and \texttt{RiskEngine}; the six specialist agents share a common \texttt{analyze} interface; \texttt{Sage} uses the \texttt{KnowledgeBase}, and the \texttt{Explainer} uses Sage's completion function and verifies its numbers. The liquidation classes create, aggregate and evaluate liquidation events.

\diagram{diagrams/04_class_diagram.png}{Class Diagram}{fig:class}{\linewidth}

\subsection{Sequence Diagrams}
This section presents four sequence diagrams for the critical workflows.

\textbf{1. Registration, login and authenticated request.} The client sends the credentials to the authentication router, which hashes or verifies the password with bcrypt, stores or reads the user, and returns a JSON Web Token. Later requests send the token as a bearer header and are authorised by decoding it (Figure~\ref{fig:seqauth}).

\diagram{diagrams/05_seq_authentication.png}{User Authentication Sequence Diagram}{fig:seqauth}{0.8\linewidth}

\textbf{2. AI signal and risk score.} For each coin the signal service fetches candles, funding and the Fear and Greed index, builds the features and asks the predictor for the ensemble probabilities. The risk engine turns the signal into a risk score, which is saved at most every thirty minutes. If any input is missing, the API returns HTTP 503 with the list of missing inputs instead of made-up values (Figure~\ref{fig:seqrisk}).

\diagram{diagrams/06_seq_signal_risk.png}{AI Signal and Risk Score Sequence Diagram}{fig:seqrisk}{\linewidth}

\textbf{3. Risk Copilot chat.} The chat router gathers live context, runs the six specialists and passes their reports to Sage. Sage retrieves sources from the knowledge base and asks Claude for a cited answer; if Claude fails or no key is configured it falls back to an OpenAI-compatible model and finally to rule-based answers (Figure~\ref{fig:seqchat}).

\diagram{diagrams/07_seq_chat.png}{Risk Copilot Chat Sequence Diagram}{fig:seqchat}{\linewidth}

\textbf{4. Liquidation intelligence and alerts.} Binance pushes forceOrder snapshots, which the stream validates, classifies and de-duplicates. Every three seconds the store loop writes new events to the database. The client requests statistics, which are computed with SQL aggregates. Every thirty seconds the alert watcher evaluates liquidation alerts with robust z-scores and a cooldown, stores notifications and sends them to Telegram (Figure~\ref{fig:seqliq}).

\diagram{diagrams/08_seq_liquidation.png}{Liquidation Intelligence and Alerts Sequence Diagram}{fig:seqliq}{\linewidth}

\subsection{State Transition Diagram}
Figure~\ref{fig:state} shows three important state machines. (a) The liquidation stream moves between connecting, connected and disconnected states; a connection that stays silent for 300 seconds is treated as broken, the gap is recorded and the stream reconnects after five seconds. (b) A liquidation alert is active until its condition is met, is then triggered and cools down before becoming active again, and can be paused by the user. (c) An AI signal is requested, its real inputs are fetched, and it becomes either available (cached for five minutes) or unavailable with the list of missing inputs.

\diagram{diagrams/09_state_diagrams.png}{State Diagrams}{fig:state}{\linewidth}

\section{Data Design}
CryptoSense stores data in a relational database accessed through SQLAlchemy \cite{sqlalchemy}. SQLite is used for local development and PostgreSQL in deployment. Figure~\ref{fig:erd} shows the entity relationship diagram. Each user owns holdings, alerts, notifications and one paper account, which owns positions and trades. Risk predictions, listings and liquidation events are independent tables. The liquidation table is indexed on event time, on (symbol, event time) and on a five-minute bucket so that window statistics and baselines can be computed with SQL aggregation. Trained models are stored as files (\texttt{xgb.joblib}, \texttt{lgbm.joblib} and \texttt{metrics.json}), and the assistant's knowledge base is stored as a TF-IDF document file.

\diagram{diagrams/10_er_diagram.png}{Entity Relationship Diagram}{fig:erd}{\linewidth}

\subsection{Data Dictionary}
{\renewcommand{\arraystretch}{1.3}
\begin{longtable}{|p{0.22\linewidth}|p{0.16\linewidth}|p{0.52\linewidth}|}
\caption{Data Dictionary for CryptoSense AI}\label{tab:datadict}\\\hline
\textbf{Entity} & \textbf{Type} & \textbf{Description}\\\hline\endfirsthead
\hline\textbf{Entity} & \textbf{Type} & \textbf{Description}\\\hline\endhead
users & Table & Account name, unique email, bcrypt password hash and creation time.\\\hline
holdings & Table & Coins owned by a user: symbol, amount, average buy price, stop-loss and take-profit percentages, and whether a limit has fired.\\\hline
alerts & Table & User alerts: symbol, kind (above, below, move, risk, liq\_large, liq\_spike, liq\_long, liq\_short, liq\_imbal), threshold value, active flag and trigger time.\\\hline
notifications & Table & Messages created when alerts or limits fire, with kind and seen flag.\\\hline
paper\_accounts & Table & Pretend cash balance and starting cash of a user's practice account.\\\hline
paper\_positions & Table & Open practice positions: symbol, amount and average price.\\\hline
paper\_trades & Table & Practice trades with side, amount, price, value, fee, realised profit and the risk and trend at the time.\\\hline
risk\_predictions & Table & History of risk scores per coin (at most one row every 30 minutes).\\\hline
listings & Table & Binance pairs with listing date, prices, 24-hour volume, volatility and risk rating.\\\hline
liquidation\_events & Table & Observed liquidations: unique event key, symbol, side, position, order type, status, quantity, filled quantity, prices, USD value and validity, event time and five-minute bucket.\\\hline
Signal & JSON & Trend, model score, votes, class probabilities, reasons, key levels and the 18 indicator values of a coin.\\\hline
Specialist report & JSON & Name, status (available, partial, unavailable), summary and data of one specialist analyst.\\\hline
\end{longtable}}

\subsection{Machine Learning Design}
Figure~\ref{fig:mlpipeline} shows the training pipeline. Hourly candles are downloaded from Binance for up to ten years, together with futures funding and the Fear and Greed history. Only hours for which funding and the Fear and Greed index were really observed are kept, so the model is never trained on filled-in values. Eighteen features (Table~\ref{tab:features}) are computed, and each hour is labelled by the price 24 hours later: above $+2.5\%$ is bullish, below $-2.5\%$ is bearish, otherwise neutral. Each coin's data is split chronologically into 70\% training, 15\% validation and 15\% test, with 24-hour gaps so that no label from a later part leaks into an earlier part. XGBoost and LightGBM are trained with class balancing and early stopping on the validation set. Their probabilities are averaged with weights equal to each model's validation macro-F1, and the ensemble is evaluated on the untouched test set against a majority-class baseline.

\diagram{diagrams/11_ml_pipeline.png}{Machine Learning Training Pipeline}{fig:mlpipeline}{\linewidth}

{\renewcommand{\arraystretch}{1.25}
\begin{longtable}{|p{0.22\linewidth}|p{0.68\linewidth}|}
\caption{Features Used by the Machine Learning Models}\label{tab:features}\\\hline
\textbf{Feature} & \textbf{Description}\\\hline\endfirsthead
\hline\textbf{Feature} & \textbf{Description}\\\hline\endhead
ret\_1, ret\_4, ret\_24 & Percentage return over the last 1, 4 and 24 hours.\\\hline
rsi\_14 & 14-period Relative Strength Index \cite{wilder1978}.\\\hline
macd\_hist, macd\_cross & MACD histogram (as \% of price) and the sign of MACD minus its signal line.\\\hline
bb\_pos, bb\_width & Position inside the Bollinger Bands and band width (\% of price).\\\hline
atr\_pct & 14-period Average True Range \cite{wilder1978} as a percentage of price.\\\hline
sma20\_dist, sma50\_dist & Distance of price from its 20- and 50-hour moving averages (\%).\\\hline
vol\_ratio & Volume relative to its 24-hour average.\\\hline
taker\_ratio & Share of volume bought aggressively over six hours (order flow).\\\hline
funding\_rate & Latest Binance futures funding rate (\%).\\\hline
fng, fng\_change\_7d & Fear and Greed index and its seven-day change.\\\hline
btc\_ret\_24, btc\_sma50\_dist & Bitcoin's 24-hour return and distance from its 50-hour average.\\\hline
\end{longtable}}

The risk score combines seven factors whose maximum points sum to 100 (Table~\ref{tab:riskfactors}). Each factor is scaled between zero and its maximum, and the explanation of the largest factor is shown to the user.

{\renewcommand{\arraystretch}{1.25}
\begin{longtable}{|p{0.22\linewidth}|p{0.1\linewidth}|p{0.56\linewidth}|}
\caption{Factors of the Explainable Risk Score}\label{tab:riskfactors}\\\hline
\textbf{Factor} & \textbf{Max} & \textbf{Basis}\\\hline\endfirsthead
\hline\textbf{Factor} & \textbf{Max} & \textbf{Basis}\\\hline\endhead
Price swings & 30 & ATR as a percentage of price (typical hourly move).\\\hline
Trend & 15 & How far price is below its 50-hour average.\\\hline
AI model view & 15 & The ensemble's bearish score (raw model score).\\\hline
Momentum stretch & 10 & How far RSI is from 50 (overbought or oversold).\\\hline
Crowded trades & 10 & Size of the futures funding rate.\\\hline
Market mood & 10 & Distance of the Fear and Greed index from neutral.\\\hline
Bitcoin trend & 10 & How far Bitcoin is below its 50-hour average.\\\hline
\end{longtable}}

\subsection{Liquidation Analytics Design}
Each forceOrder message is validated: the symbol must exist, the side must be BUY or SELL and the price and quantity must be positive. A forced SELL closes a long position and is counted as a \textit{long liquidation}; a forced BUY closes a short position and is counted as a \textit{short liquidation}. The USD value is the filled quantity multiplied by the average fill price, and is marked valid only for USD-M contracts quoted in a stablecoin. Duplicate events are removed with a key built from the symbol, time, side, quantity and price.

Spikes are detected by comparing the last five minutes with a rolling baseline of five-minute buckets from the previous 24 hours. The robust z-score is
\begin{equation}
z = \frac{x - \mathrm{median}(B)}{\max\left(1.4826 \cdot \mathrm{MAD}(B),\ m/4\right)}
\end{equation}
where $x$ is the current value, $B$ is the baseline, MAD is the median absolute deviation \cite{rousseeuw1993} and $m$ is the minimum notable amount (\$250{,}000 by default). The floor $m/4$ prevents a nearly silent baseline from turning a small amount into a very large score. A spike requires both $z \ge 4$ and $x \ge m$. Bursts (at least 15 events for one coin within 60 seconds) and shifts in the long share of liquidations (at least 35 percentage points) are also flagged. At least one hour of history is required before spikes are judged.

\subsection{Grounded Explanation Design}
Every explanation starts from facts computed by the analytics and a rules-based answer to each question. If a language model is configured, it receives only the facts and the questions, with instructions to use no other numbers. The explainer then extracts every number from the model's text and checks that it matches a number in the facts at the precision it was written (allowing rounding, ``K/M/B'' suffixes and percentages). A few small numbers that describe time windows or counts (0 to 10, 15, 24, 30, 60 and 100) are always allowed. If any number cannot be traced, the model's answer is discarded and the rules-based answer is shown instead.

\section{Human Interface Design}
The interface uses a dark theme with a gold accent, a left navigation sidebar and a top bar with live BTC and ETH prices and a Live/Offline badge. Every panel has a short subtitle and, where useful, a lightbulb button that opens a plain-language explanation from Sage. The screens below were captured from the running application with live data.

\subsection{Screen Images}
\diagram{screens/overview.png}{Overview Dashboard with Sage's Market Read and Risk Score}{fig:scr-overview}{\linewidth}
\diagram{screens/copilot_answer.png}{CryptoSentinel Risk Copilot Answering with Six Specialist Agents}{fig:scr-copilot}{\linewidth}
\diagram{screens/liquidations.png}{Liquidation Intelligence Screen}{fig:scr-liq}{\linewidth}
\diagram{screens/news.png}{News Screen with Sentiment Tags}{fig:scr-news}{\linewidth}
\diagram{screens/alerts.png}{Alerts Screen}{fig:scr-alerts}{\linewidth}
\diagram{screens/portfolio.png}{Portfolio Screen}{fig:scr-portfolio}{\linewidth}
\diagram{screens/whales.png}{Whale Watch and Smart-Money Screen}{fig:scr-whales}{\linewidth}
\diagram{screens/listings.png}{New Listings Screen}{fig:scr-listings}{\linewidth}
\diagram{screens/academy.png}{Academy Screen}{fig:scr-academy}{\linewidth}

\subsection{Screen Objects and Actions}
{\renewcommand{\arraystretch}{1.3}
\begin{longtable}{|p{0.24\linewidth}|p{0.24\linewidth}|p{0.42\linewidth}|}
\caption{Screen Objects and Their Actions}\label{tab:screenobjects}\\\hline
\textbf{Screen Object} & \textbf{Action} & \textbf{Description}\\\hline\endfirsthead
\hline\textbf{Screen Object} & \textbf{Action} & \textbf{Description}\\\hline\endhead
Coin chips (Overview) & Select coin & Switches the market read, chart and risk score to the chosen coin.\\\hline
Interval pills (chart) & Change interval & Reloads real candles for 15m, 1h, 4h or 1d.\\\hline
Lightbulb button & Explain & Opens Sage's plain-language explanation of the panel.\\\hline
Volatility what-if & Simulate & Recomputes the risk score at 1.0x, 1.5x or 2.0x volatility.\\\hline
Open Copilot / Send & Ask Sage & Sends the question to the six-specialist chat and shows the answer, specialists and sources.\\\hline
Window and coin pills (Liquidations) & Filter & Recomputes liquidation statistics for 5m--24h and BTC, ETH, SOL, BNB, XRP or other coins.\\\hline
Explain button (Liquidations) & Explain & Returns verified answers to six questions about liquidation activity.\\\hline
Create alert & Create & Opens the alert form for price, move, risk and liquidation alerts.\\\hline
Add holding & Create & Adds a coin, amount and buy price to the portfolio.\\\hline
Limit sliders & Set limits & Sets stop-loss and take-profit percentages.\\\hline
Buy / Sell (Paper trading) & Trade & Executes a practice trade at the live price with a 0.1\% fee.\\\hline
Sign in & Authenticate & Opens login or registration.\\\hline
\end{longtable}}

{\renewcommand{\arraystretch}{1.3}
\begin{longtable}{|p{0.24\linewidth}|p{0.14\linewidth}|p{0.52\linewidth}|}
\caption{Tools and Technologies for CryptoSense AI}\label{tab:tools}\\\hline
\textbf{Tool / Library} & \textbf{Version} & \textbf{Rationale}\\\hline\endfirsthead
\hline\textbf{Tool / Library} & \textbf{Version} & \textbf{Rationale}\\\hline\endhead
Python & 3.11 & Backend language with a rich data-science ecosystem.\\\hline
FastAPI, Uvicorn & 0.110+ & Fast asynchronous REST API with automatic OpenAPI documentation.\\\hline
SQLAlchemy & 2.x & Database access for SQLite and PostgreSQL.\\\hline
XGBoost, LightGBM & 2.x, 4.x & Gradient-boosted tree models for tabular features.\\\hline
scikit-learn, pandas, NumPy & 1.4+, 2.x, 1.26+ & Metrics, TF-IDF search and data processing \cite{pedregosa2011sklearn}.\\\hline
vaderSentiment & 3.3 & Transparent news sentiment scoring.\\\hline
websockets, httpx & 12+, 0.27+ & Binance WebSocket streams and HTTP clients.\\\hline
bcrypt, PyJWT & 4.x, 2.x & Password hashing and JSON Web Tokens.\\\hline
Anthropic SDK & 2.13+ & Claude language model access for explanations.\\\hline
Flutter, Dart & 3.35 & Cross-platform web user interface.\\\hline
fl\_chart & 1.3 & Interactive bar charts for liquidation analytics.\\\hline
draw.io & -- & Design diagrams.\\\hline
\end{longtable}}

\section{External APIs/SDKs}
{\renewcommand{\arraystretch}{1.3}
\begin{longtable}{|p{0.24\linewidth}|p{0.26\linewidth}|p{0.4\linewidth}|}
\caption{External APIs and SDKs}\label{tab:apis}\\\hline
\textbf{API / SDK} & \textbf{Purpose} & \textbf{Integration Details}\\\hline\endfirsthead
\hline\textbf{API / SDK} & \textbf{Purpose} & \textbf{Integration Details}\\\hline\endhead
Binance Spot REST and WebSocket & Prices, candles, tickers, listings & Public endpoints, no key \cite{binanceapi}; candles and tickers via REST, live tickers via WSS.\\\hline
Binance USD-M Futures & Funding, open interest, futures candles, liquidations & Public REST endpoints and the all-market forceOrder WebSocket stream \cite{binancefutures}.\\\hline
alternative.me & Fear and Greed index & Public JSON API (daily values) \cite{feargreed}.\\\hline
CoinDesk / Cointelegraph RSS, CryptoPanic & News headlines & RSS by default; CryptoPanic when a key is set.\\\hline
Whale Alert & Large on-chain transfers & Optional key; transfers above \$1 million in the last hour.\\\hline
Quiver Quantitative & Congress and insider trades & Optional key; filtered to crypto-linked stocks.\\\hline
Anthropic Claude & Answer phrasing & Optional key; used first by Sage and the explainer \cite{anthropic}.\\\hline
Groq (OpenAI-compatible) & Fallback answer phrasing & Optional key; any OpenAI-compatible endpoint can be configured.\\\hline
TypeSafe Jev & News judgement & Optional key; judges headline sentiment and impact.\\\hline
Telegram Bot API & Alert delivery & Optional bot token and chat ID \cite{telegram}.\\\hline
\end{longtable}}

\subsection{REST API Endpoints}
{\renewcommand{\arraystretch}{1.2}\small
\begin{longtable}{|p{0.1\linewidth}|p{0.38\linewidth}|p{0.42\linewidth}|}
\caption{REST API Endpoints of the Backend}\label{tab:endpoints}\\\hline
\textbf{Method} & \textbf{Path} & \textbf{Purpose}\\\hline\endfirsthead
\hline\textbf{Method} & \textbf{Path} & \textbf{Purpose}\\\hline\endhead
POST & /api/auth/register, /login & Create an account; log in and receive a token.\\\hline
GET & /api/auth/me & Current user.\\\hline
GET & /api/market/tickers, /klines, /sentiment & Prices, candles, Fear and Greed with funding and open interest.\\\hline
GET & /api/signals, /api/signals/\{symbol\} & AI signals (503 with missing inputs if unavailable).\\\hline
GET, POST & /api/signals/model-info, /reload & Model metrics; reload models after training.\\\hline
GET & /api/risk, /api/risk/\{symbol\}, /portfolio & Risk scores and portfolio risk.\\\hline
GET, POST & /api/risk/\{symbol\}/copilot, /simulate, /history & Votes, twins, levels; what-if; history.\\\hline
GET, POST, DELETE, PUT & /api/portfolio, /\{symbol\}, /\{symbol\}/limits & Holdings and profit/loss limits.\\\hline
GET, POST & /api/paper, /trade, /trades, /reset & Paper-trading account.\\\hline
GET, POST, PATCH, DELETE & /api/alerts, /\{id\}/toggle, /\{id\} & Alerts.\\\hline
GET, POST & /api/notifications, /seen & Notifications.\\\hline
GET & /api/news, /news/compare, /whales, /quiver, /listings & Discovery feeds.\\\hline
GET & /api/liquidations & Recent one-hour liquidation snapshot.\\\hline
GET & /api/liquidations/status, /stats, /timeline, /history & Stream health, window statistics, chart data, history.\\\hline
GET & /api/liquidations/anomalies, /context/\{symbol\}, /explain & Spike detection, market context, verified explanation.\\\hline
POST, GET & /api/chat, /agents, /knowledge, /knowledge/rebuild & Copilot chat and knowledge base.\\\hline
\end{longtable}}

\section{Deployment}
Figure~\ref{fig:deployment} shows the deployment. The Flutter application is built for the web and runs in the user's browser. The backend runs on Railway (or a local computer) as a Uvicorn ASGI server hosting the FastAPI application, its background tasks, the trained model files and the knowledge store. A Railway PostgreSQL database stores the application data; locally, SQLite is used. Configuration such as the database URL, the JWT secret and optional API keys is read from environment variables. For local use, \texttt{SETUP.bat} installs the dependencies, \texttt{TRAIN\_AI.bat} trains the models and \texttt{START.bat} starts the backend and then the frontend in Chrome.

\diagram{diagrams/12_deployment_diagram.png}{Deployment Diagram}{fig:deployment}{\linewidth}

\section{Chapter Summary}
This chapter presented the design of CryptoSense AI: its modules, its layered, event-driven client-server architecture, the activity, use case, class, sequence and state diagrams, the data, machine learning, liquidation and grounded explanation designs, the human interface, the tools and external APIs, the REST endpoints and the deployment. The next chapter describes how the system was tested and evaluated.

\fypchapter{Testing \& Evaluation}{ch:testing}

\chapteroverview
This chapter presents the testing and evaluation of CryptoSense AI. The backend was verified with an automated suite of 65 unit and integration tests, the frontend with static analysis and a production web build, and the running system with live end-to-end checks against real Binance data. The machine learning models were evaluated on unseen, chronologically later data and compared with a majority-class baseline. All results reported here come from tests and measurements that were actually executed.

\section{Testing Approach and Environment}
\begin{itemize}
\item \textbf{Automated backend tests:} 65 tests written with Python's \texttt{unittest} framework, run with \texttt{python -m unittest discover -s tests}. External services are replaced by test doubles, and an in-memory SQLite database is used for storage tests, so the tests are fast and repeatable. Result of the final run: \textbf{65 tests, all passed}.
\item \textbf{Frontend checks:} \texttt{flutter analyze} reported no errors or warnings in the application code (only style suggestions), and \texttt{flutter build web} compiled the application successfully.
\item \textbf{Live end-to-end checks:} the backend and frontend were run together on a Windows PC (Python 3.11.9, Flutter 3.35.4, Chrome) against live Binance data, and every public GET endpoint and the main screens were exercised.
\item \textbf{Model evaluation:} chronological train/validation/test splits with gaps, a majority-class baseline, and a comparison of models on the same unseen period.
\end{itemize}

\section{Unit Testing}
Unit tests verify individual components in isolation. Each case below corresponds to one or more automated tests in the \texttt{backend/tests} folder.

\testcase{Unit Testing 1: Liquidation Event Parsing and Classification}
{Verify that forced SELL orders are classified as long liquidations, forced BUY orders as short liquidations, and that the USD value equals filled quantity times average price.}
{Correct position, side and USD value; USD marked invalid for coin-margined or non-stablecoin contracts.}
{All parsing tests passed (test\_sell\_is\_long\_and\_buy\_is\_short, test\_usd\_uses\_filled\_quantity\_times\_average\_price, test\_usd\_not\_valid\_for\_coin\_margined\_or\_non\_stable\_quote).}
{Pass}

\testcase{Unit Testing 2: Invalid, Missing and Duplicate Data}
{Ensure that malformed messages, missing fields, zero prices and duplicates are rejected.}
{Invalid records produce no events; a repeated event is stored only once.}
{Verified by test\_invalid\_and\_missing\_data\_is\_skipped, test\_duplicates\_are\_dropped and test\_store\_skips\_existing\_rows.}
{Pass}

\testcase{Unit Testing 3: Time-Window Statistics}
{Validate totals, long/short ratio, largest event, coin filters and zero-filled, UTC-aligned timelines for 5m to 24h windows.}
{Totals cover the whole window (not only the returned events); timelines are sorted and aligned to bucket boundaries.}
{Verified by test\_window\_totals\_ratio\_and\_largest, test\_symbol\_and\_other\_filters, test\_timeline\_is\_zero\_filled\_and\_utc\_aligned and test\_totals\_cover\_whole\_window\_not\_just\_limit.}
{Pass}

\testcase{Unit Testing 4: Liquidation Spike Detection}
{Check that spikes are detected only above the robust z-score and minimum amount, and that normal activity and small amounts on a quiet baseline are not flagged.}
{Long spike, total spike and imbalance shift detected for injected test data; no anomaly for normal data; ``insufficient baseline'' with less than one hour of history.}
{Verified by test\_spike\_detected\_only\_above\_z\_and\_minimum, test\_small\_amount\_on\_quiet\_baseline\_is\_not\_a\_spike, test\_normal\_activity\_is\_quiet, test\_burst and test\_db\_detection\_needs\_baseline\_then\_finds\_spike.}
{Pass}

\testcase{Unit Testing 5: Grounded Explanation (Number Verification)}
{Ensure that numbers written by a language model are accepted only if they appear in the computed facts.}
{Rounded values and K/M suffixes are accepted; invented prices and probabilities are rejected and the rules-based answer is returned.}
{Verified by test\_grounded\_numbers\_pass\_including\_rounding\_and\_suffixes, test\_invented\_numbers\_are\_caught and test\_llm\_answer\_with\_invented\_numbers\_falls\_back\_to\_rules.}
{Pass}

\testcase{Unit Testing 6: No Fabricated Data}
{Verify that missing live inputs produce ``unavailable'' instead of made-up values.}
{Empty prices when Binance fails; signals unavailable (HTTP 503 with the list of missing inputs) when candles, funding, the Fear and Greed index or a trained model are missing; whales and news unavailable without a source.}
{Verified by the six tests in test\_no\_fake\_data.py.}
{Pass}

\testcase{Unit Testing 7: Specialist Agents}
{Check that the six specialists report their status honestly and that the orchestrator returns all six in order.}
{Partial status when live figures are missing; unavailable whale or news feeds are not used as live evidence.}
{Verified by the six tests in test\_agents.py.}
{Pass}

\testcase{Unit Testing 8: Language Model Fallback Chain}
{Verify that Claude is used first, that failures fall back to the OpenAI-compatible model and then to rules, and that a refusal triggers the fallback.}
{Correct provider order and request shape.}
{Verified by the four Sage tests in test\_integrations.py.}
{Pass}

\section{Functional Testing}
Functional testing validates that each module performs its intended operations in the running system with live data.

\testcase{Functional Testing 1: Live Market Module}
{Confirm that real prices stream into the interface and that the chart loads real candles.}
{Live badge shown; BTC and ETH prices update; chart shows Binance candles.}
{Live prices and candles shown (Figure~\ref{fig:scr-overview}); with Binance unreachable the screen states that live data is unavailable.}
{Pass}

\testcase{Functional Testing 2: AI Signals and Risk Score}
{Verify that signals and risk scores are produced from live data and explained.}
{Signals for all six coins with model votes; risk score with seven factors.}
{All six signals returned by the ensemble; BTC risk shown with its factors (36\%, moderate, at the time of testing).}
{Pass}

\testcase{Functional Testing 3: Risk Copilot Chat}
{Verify that a question is answered by the backend with the six specialists and sources.}
{Answer with specialist cards and cited sources.}
{The question ``How risky is BTC right now?'' was answered in about 2.4 seconds with six specialist reports and three sources (Figure~\ref{fig:scr-copilot}).}
{Pass}

\testcase{Functional Testing 4: Liquidation Intelligence}
{Confirm that the liquidation stream receives real events and that the screen shows statistics and charts.}
{Connected stream, valid events, totals and charts.}
{In a live run the stream received 94 messages with 0 invalid and 0 duplicates; the screen showed totals, ratio, largest event, most affected coin and charts (Figure~\ref{fig:scr-liq}).}
{Pass}

\testcase{Functional Testing 5: Public API Endpoints}
{Verify that every public GET endpoint responds.}
{HTTP 200 with real or clearly ``unavailable'' data.}
{All 19 public GET endpoints returned HTTP 200; whale and Quiver feeds correctly reported ``unavailable'' because no keys were configured.}
{Pass}

\section{Business Rules Testing}
This testing ensures that the internal rules of the system behave correctly.

{\renewcommand{\arraystretch}{1.3}
\begin{longtable}{|p{0.26\linewidth}|p{0.32\linewidth}|p{0.32\linewidth}|}
\caption{AI Signal Availability Business Rule}\label{tab:br-signal}\\\hline
\textbf{Condition} & \textbf{Description} & \textbf{Expected Action}\\\hline
All inputs real & Candles, funding and Fear and Greed fetched; trained model present & Compute and return the signal\\\hline
Any input missing & One or more sources unreachable & Return HTTP 503 with the missing inputs; never substitute values\\\hline
Low confidence & Top directional score below 0.55 & Show the signal as neutral\\\hline
\end{longtable}}

{\renewcommand{\arraystretch}{1.3}
\begin{longtable}{|p{0.26\linewidth}|p{0.32\linewidth}|p{0.32\linewidth}|}
\caption{Liquidation Alert Business Rule}\label{tab:br-liq}\\\hline
\textbf{Condition} & \textbf{Description} & \textbf{Expected Action}\\\hline
Condition met & Event size or robust z-score reaches the user's threshold and the minimum amount & Create a notification and send it to Telegram\\\hline
Cooling down & Alert fired within the cooldown period (30 minutes) & Do not notify again\\\hline
Insufficient history & Less than one hour of stored history & Do not judge spikes\\\hline
\end{longtable}}

{\renewcommand{\arraystretch}{1.3}
\begin{longtable}{|p{0.26\linewidth}|p{0.32\linewidth}|p{0.32\linewidth}|}
\caption{Grounded Explanation Business Rule}\label{tab:br-explain}\\\hline
\textbf{Condition} & \textbf{Description} & \textbf{Expected Action}\\\hline
No language model configured & No API key & Return the rules-based answer\\\hline
All numbers traceable & Every number matches the facts & Return the model's answer, marked as verified\\\hline
Untraceable number & A number not found in the facts & Discard the model's answer and return the rules-based answer\\\hline
\end{longtable}}

\section{Integration Testing}
\testcase{Integration Testing 1: Chat Endpoint with Specialists}
{Ensure that the chat endpoint gathers context, runs all six specialists and returns them with the answer.}
{Response contains the answer and six specialist reports.}
{Verified by test\_chat\_returns\_all\_six\_specialist\_reports and in the running system.}
{Pass}

\testcase{Integration Testing 2: Liquidation Store, Statistics and Explanation}
{Validate the flow from stored events to the statistics, timeline, history, anomaly and explanation endpoints.}
{Correct totals and six rules-based answers without a language model.}
{Verified by test\_stats\_timeline\_history\_status and test\_explain\_uses\_rules\_without\_llm.}
{Pass}

\testcase{Integration Testing 3: Stream Disconnection and Reconnection}
{Verify that the stream reconnects after a network error and records the gap.}
{One reconnect counted, gap closed, events received after reconnecting.}
{Verified by test\_disconnect\_then\_reconnect\_is\_tracked.}
{Pass}

\testcase{Integration Testing 4: Alerts and Telegram}
{Verify that alerts are evaluated with the cooldown and delivered to Telegram only when configured.}
{One notification per event; no message without credentials.}
{Verified by test\_evaluate\_fires\_once\_then\_cools\_down, the Telegram tests and the watcher tests.}
{Pass}

\testcase{Integration Testing 5: Frontend and Backend}
{Ensure that the web application consumes the backend correctly.}
{All screens load data from the backend without errors.}
{During the live run every request from the browser returned HTTP 200 and no errors were logged by the backend or the application.}
{Pass}

\section{Machine Learning Evaluation}
\subsection{Deployed Ensemble}
The deployed models were trained on two years of hourly data (24-hour horizon, $\pm2.5\%$ threshold) with 72{,}984 training, 15{,}414 validation and 15{,}708 test rows. Table~\ref{tab:ml-deployed} shows the results on the test set. ``Confident accuracy'' is the accuracy on the predictions whose top score is at least 0.5, and ``coverage'' is the share of predictions that are this confident.

{\renewcommand{\arraystretch}{1.3}
\begin{longtable}{|p{0.26\linewidth}|c|c|c|c|}
\caption{Test-Set Results of the Deployed Models}\label{tab:ml-deployed}\\\hline
\textbf{Model} & \textbf{Accuracy} & \textbf{Macro-F1} & \textbf{Confident acc.} & \textbf{Coverage}\\\hline
XGBoost & 0.556 & 0.414 & 0.782 & 0.351\\\hline
LightGBM & 0.544 & 0.397 & 0.656 & 0.517\\\hline
Ensemble & 0.555 & 0.405 & 0.726 & 0.419\\\hline
Always ``neutral'' (baseline) & 0.712 & 0.277 & -- & --\\\hline
\end{longtable}}

The models are trained with class balancing so that they try to recognise the rarer bullish and bearish periods; this improves macro-F1 (0.405 versus 0.277 for the baseline) but lowers plain accuracy, because about 70\% of hours are neutral. Accuracy alone is therefore not a good measure of skill for this task. When the models are confident, their accuracy rises to 0.73--0.78, which is why low-confidence directional calls are shown as neutral.

\subsection{Effect of Training History}
Models were retrained on five and ten years of requested history and compared with the deployed two-year models on the \textbf{same} final 100 days (14{,}256 hourly rows), which none of the three models saw during training. For the ten-year run, only hours with real funding and Fear and Greed data were used, so the training data actually starts between September 2019 (BTC) and September 2020 (SOL).

{\renewcommand{\arraystretch}{1.3}
\begin{longtable}{|p{0.28\linewidth}|c|c|c|c|c|}
\caption{Comparison on the Same Unseen 100 Days}\label{tab:ml-history}\\\hline
\textbf{History} & \textbf{Accuracy} & \textbf{Macro-F1} & \textbf{Log loss} & \textbf{Fall Brier} & \textbf{Fall AUC}\\\hline
2 years (deployed) & 0.518 & 0.395 & 0.987 & 0.160 & 0.573\\\hline
5 years & 0.564 & 0.407 & 0.973 & 0.160 & 0.613\\\hline
10 years requested & 0.597 & 0.407 & 0.924 & 0.152 & 0.608\\\hline
Always ``neutral'' & 0.687 & -- & -- & -- & --\\\hline
\end{longtable}}

Longer history clearly helps: the ten-year model has the best accuracy, log loss and Brier score, and both longer histories identify falls better (ROC-AUC about 0.61 instead of 0.57). However, all models remain weak predictors of direction, which is consistent with the efficient market hypothesis \cite{fama1970}.

\subsection{Calibration and Downside-Risk Experiment}
A further experiment tested two improvements on the most recent year of unseen data (25 October 2025 to 22 November 2026, 53{,}574 rows): twelve additional volatility and calendar features computed only from past data, and probability calibration (multinomial Platt scaling fitted on a held-out part of the validation set). It also tested a different target: whether the price \textit{touches} 2.5\% below the current close at any time in the next 24 hours, which measures downside risk rather than direction.

{\renewcommand{\arraystretch}{1.3}
\begin{longtable}{|p{0.42\linewidth}|c|c|c|c|}
\caption{Calibration and Downside-Risk Experiment (Unseen Year)}\label{tab:ml-experiment}\\\hline
\textbf{Model} & \textbf{Accuracy} & \textbf{Log loss} & \textbf{Brier} & \textbf{AUC}\\\hline
\multicolumn{5}{|l|}{\textit{Direction (3 classes), Brier and AUC for the bearish class}}\\\hline
Baseline (class frequencies) & 0.689 & 0.867 & 0.146 & 0.500\\\hline
LightGBM, current features & 0.687 & 0.811 & 0.142 & 0.626\\\hline
LightGBM + volatility features, calibrated & 0.689 & 0.796 & 0.140 & 0.643\\\hline
\multicolumn{5}{|l|}{\textit{Downside touch of $-2.5\%$ within 24 hours (base rate 0.354)}}\\\hline
Baseline (base rate) & 0.646 & 0.671 & 0.239 & 0.500\\\hline
Logistic regression on ATR only & 0.662 & 0.620 & 0.215 & 0.670\\\hline
LightGBM + volatility features & 0.668 & 0.598 & 0.207 & 0.689\\\hline
\end{longtable}}

Calibration and volatility features reduce log loss and bring the Brier score for falls to 0.140, better than the baseline of 0.146, while discrimination improves to an AUC of 0.643. The downside-touch target is more predictable than direction (AUC 0.689), largely because volatility is persistent: a model using ATR alone already reaches an AUC of 0.670. These results answer research questions RQ1--RQ3: public market data contains a weak but real signal, more history helps, and downside risk is more predictable than direction. They also show that much higher figures (for example an AUC of 0.80 for 24-hour direction) are not achievable from these data without information leakage. Further gains require new information, such as options-implied volatility and recorded open interest and liquidation history.

\section{Chapter Summary}
This chapter described the testing and evaluation of CryptoSense AI. All 65 automated backend tests passed, the frontend compiled without errors, and live end-to-end checks confirmed that the system works with real Binance data and reports unavailable inputs honestly. The machine learning models were evaluated against a baseline on unseen data: they provide a weak but measurable signal, which improves with longer history, calibration and volatility features, and downside risk proved more predictable than direction. The next chapter concludes the project and proposes future work.

\fypchapter{Conclusion and Future Work}{ch:conclusion}

\chapteroverview
This chapter presents the overall conclusion of the CryptoSense AI project. It summarises the contributions and findings of the system and proposes directions for future improvement and research.

\section{Conclusion}
CryptoSense AI was developed to help everyday cryptocurrency users understand market conditions and risk in plain language. The system brings together live prices, futures funding and open interest, liquidations, news sentiment, whale and smart-money activity, and machine learning signals in a single web application, and explains them through an explainable risk score, a Risk Copilot and a conversational assistant supported by six specialist analysts.

The project achieved its objectives. Live market data is streamed from Binance and never replaced by fabricated values (PO1). An XGBoost and LightGBM ensemble was trained and honestly evaluated against a baseline with a leak-free chronological split (PO2). The seven-factor risk score, model votes, historical twins, key levels and what-if simulation make risk explainable (PO3). Sage coordinates six specialists, retrieves cited sources and uses language models only to phrase verified results (PO4). The Liquidation Intelligence module records observed liquidations, computes multi-window statistics, detects unusual activity with robust statistics and raises alerts with a cooldown (PO5). Portfolio limits, alerts with Telegram delivery and paper trading help users manage risk (PO6).

The evaluation showed that public market data contains a weak but measurable signal: longer training history, volatility features and calibration improve the models, and downside risk is more predictable than direction. Equally important, the project showed how to present such models responsibly: by stating their limits, comparing them with a baseline, labelling raw scores as scores, and refusing to show made-up data or unverified numbers.

\section{Future Work}
Several enhancements can further improve accuracy, coverage and usefulness:
\begin{itemize}
\item \textbf{Calibrated downside-risk models:} train separate, calibrated models for precisely defined downside events at several horizons (for example a fall of at least 1\% within one hour, 2\% within four hours and 3\% within 24 hours), and show a horizon only after it passes out-of-sample validation (Brier skill above the baseline and a low calibration error).
\item \textbf{New information sources:} add options-implied volatility (Deribit DVOL), recorded open-interest history, order-book depth and the liquidation history that CryptoSense now records, which are expected to improve risk prediction.
\item \textbf{Risk Intelligence Specialist:} separate downside probability, volatility risk, drawdown severity, market stress, liquidity risk and derivatives risk, and validate the 0--100 score boundaries against realised outcomes.
\item \textbf{Support and Resistance Engine:} detect price zones from swing points across five timeframes, score their strength, recognise breakouts, rejections and false breakouts, and draw the zones on the chart.
\item \textbf{Market Regime and Anomaly Detection:} classify trending, ranging and high-volatility regimes, and extend the robust anomaly detection to price, volume, volatility, funding, open interest, order-book imbalance, stablecoin deviations and sentiment.
\item \textbf{Native whale tracking:} track large Bitcoin and Ethereum transfers from public blockchain data, with clear separation between confirmed and uncertain exchange labels.
\item \textbf{Main AI orchestrator and monitoring:} coordinate all modules with data-freshness checks, timeouts, request tracing, model version tracking and module health monitoring.
\item \textbf{Live forward testing:} store every live prediction, score it once its horizon has passed, and show real live accuracy on the dashboard, with alerts when performance degrades.
\item \textbf{Mobile applications:} publish Android and iOS builds from the same Flutter code base.
\end{itemize}

\section{Chapter Summary}
This chapter concluded the CryptoSense AI project. The system meets its objectives by providing a live, explainable and honest view of cryptocurrency market risk, combining machine learning, transparent risk scoring, liquidation intelligence and a grounded conversational assistant. Future work focuses on calibrated downside-risk models, new data sources, additional analytical engines, an orchestrator with monitoring, live forward testing and mobile applications.


% ======================================================= references
\clearpage
\thispagestyle{empty}
\vspace*{\fill}
\begin{center}{\LARGE\bfseries References}\end{center}
\vspace*{\fill}
\clearpage
\addcontentsline{toc}{chapter}{References}
% The References title page above replaces the heading thebibliography would add.
\makeatletter
\let\origchapter\chapter
\renewcommand{\chapter}{\@ifstar{\@gobble}{\origchapter}}
\makeatother
\begin{thebibliography}{99}

\bibitem{chen2016xgboost}
T. Chen and C. Guestrin, ``XGBoost: A scalable tree boosting system,'' in \textit{Proceedings of the 22nd ACM SIGKDD International Conference on Knowledge Discovery and Data Mining}, 2016, pp. 785--794.

\bibitem{ke2017lightgbm}
G. Ke, Q. Meng, T. Finley, T. Wang, W. Chen, W. Ma, Q. Ye, and T.-Y. Liu, ``LightGBM: A highly efficient gradient boosting decision tree,'' in \textit{Advances in Neural Information Processing Systems 30}, 2017, pp. 3146--3154.

\bibitem{pedregosa2011sklearn}
F. Pedregosa \textit{et al.}, ``Scikit-learn: Machine learning in Python,'' \textit{Journal of Machine Learning Research}, vol. 12, pp. 2825--2830, 2011.

\bibitem{hutto2014vader}
C. J. Hutto and E. Gilbert, ``VADER: A parsimonious rule-based model for sentiment analysis of social media text,'' in \textit{Proceedings of the 8th International AAAI Conference on Weblogs and Social Media (ICWSM)}, 2014.

\bibitem{lewis2020rag}
P. Lewis \textit{et al.}, ``Retrieval-augmented generation for knowledge-intensive NLP tasks,'' in \textit{Advances in Neural Information Processing Systems 33}, 2020.

\bibitem{mcnally2018}
S. McNally, J. Roche, and S. Caton, ``Predicting the price of Bitcoin using machine learning,'' in \textit{Proceedings of the 26th Euromicro International Conference on Parallel, Distributed and Network-based Processing (PDP)}, 2018, pp. 339--343.

\bibitem{chen2020bitcoin}
Z. Chen, C. Li, and W. Sun, ``Bitcoin price prediction using machine learning: An approach to sample dimension engineering,'' \textit{Journal of Computational and Applied Mathematics}, vol. 365, 112395, 2020.

\bibitem{fama1970}
E. F. Fama, ``Efficient capital markets: A review of theory and empirical work,'' \textit{The Journal of Finance}, vol. 25, no. 2, pp. 383--417, 1970.

\bibitem{lopezdeprado2018}
M. L\'opez de Prado, \textit{Advances in Financial Machine Learning}. Hoboken, NJ, USA: Wiley, 2018.

\bibitem{brier1950}
G. W. Brier, ``Verification of forecasts expressed in terms of probability,'' \textit{Monthly Weather Review}, vol. 78, no. 1, pp. 1--3, 1950.

\bibitem{niculescu2005}
A. Niculescu-Mizil and R. Caruana, ``Predicting good probabilities with supervised learning,'' in \textit{Proceedings of the 22nd International Conference on Machine Learning (ICML)}, 2005, pp. 625--632.

\bibitem{rousseeuw1993}
P. J. Rousseeuw and C. Croux, ``Alternatives to the median absolute deviation,'' \textit{Journal of the American Statistical Association}, vol. 88, no. 424, pp. 1273--1283, 1993.

\bibitem{wilder1978}
J. W. Wilder, \textit{New Concepts in Technical Trading Systems}. Greensboro, NC, USA: Trend Research, 1978.

\bibitem{binanceapi}
Binance, ``Binance Spot API documentation,'' Available: \url{https://developers.binance.com/docs/binance-spot-api-docs}, [Accessed: November 2026].

\bibitem{binancefutures}
Binance, ``USD-M Futures: All market liquidation order streams,'' Available: \url{https://developers.binance.com/docs/derivatives/usds-margined-futures/websocket-market-streams/All-Market-Liquidation-Order-Streams}, [Accessed: November 2026].

\bibitem{fastapi}
FastAPI, ``FastAPI documentation,'' Available: \url{https://fastapi.tiangolo.com}, [Accessed: November 2026].

\bibitem{flutter}
Flutter by Google, ``Flutter documentation,'' Available: \url{https://docs.flutter.dev}, [Accessed: November 2026].

\bibitem{sqlalchemy}
SQLAlchemy, ``SQLAlchemy 2.0 documentation,'' Available: \url{https://docs.sqlalchemy.org}, [Accessed: November 2026].

\bibitem{feargreed}
Alternative.me, ``Crypto Fear \& Greed Index API,'' Available: \url{https://alternative.me/crypto/fear-and-greed-index/}, [Accessed: November 2026].

\bibitem{anthropic}
Anthropic, ``Claude API documentation,'' Available: \url{https://docs.anthropic.com}, [Accessed: November 2026].

\bibitem{telegram}
Telegram, ``Telegram Bot API,'' Available: \url{https://core.telegram.org/bots/api}, [Accessed: November 2026].

\end{thebibliography}
\let\chapter\origchapter


\end{document}
